% !TeX encoding = UTF-8
% !TeX spellcheck = pl_PL

\def\filename{podsumowanie}

\chapter{Podsumowanie}\label{ch:podsumowanie}
W niniejszej pracy podjęto tematykę dynamicznie rozwijających się technologii sztucznej inteligencji w kontekście diagnostyki zmian skórnych.
Przedstawiono potencjał konwolucyjnych sieci neuronowych w obszarze \textit{computer-aided diagnosis}, jednocześnie zwracając uwagę na istniejące ograniczenia, takie jak ramy etyczne i regulacyjne, kwestie poufności oraz bezpieczeństwa danych pacjentów, brak wyjaśnialności decyzji oraz brak niezależnych i reprezentatywnych zbiorów danych.
Zauważono również rosnącą rolę urządzeń mobilnych jako narzędzi wspierających pacjentów w monitorowaniu zmian skórnych oraz zwiększaniu świadomości dotyczącej wczesnej diagnostyki.

W pracy przeanalizowano aktualną literaturę dotyczącą zastosowania konwolucyjnych sieci neuronowych w medycynie, ze szczególnym uwzględnieniem diagnostyki dermatologicznej, oraz omówiono istniejące rozwiązania mobilne w tej dziedzinie.
Zaproponowano autorski system obejmujący model \textit{pre-trained} przekonwertowany do formatu TensorFlow Lite oraz aplikację mobilną na system Android, wykonującą inferencję bezpośrednio na urządzeniu pacjenta.
Przedstawiono dostępne repozytoria i zbiory zmian skórnych, ich charakterystykę oraz luki, które wypełniają w kontekście badań nad sztuczną inteligencją w dermatologii.

Po przeprowadzeniu eksperymentów na różnych zestawach danych ostatecznie zdecydowano się wykorzystać zbiór PAD-UFES-20, zawierający zdjęcia wykonane aparatem telefonu komórkowego, co stanowiło przedmiot dalszych badań i zapewniło wysoką skuteczność pomimo zmienności obrazów.
W procesie trenowania zastosowano technikę \textit{transfer learning}, polegającą na dostosowaniu wcześniej wytrenowanego modelu do specyfiki nowych danych oraz zestaw metod augmentacji obrazu, w tym mixup, zwiększających zdolność uogólniania modeli.
Zbadano skuteczność siedmiu wersji architektury konwolucyjnej.
Najlepsze wyniki osiągnął model ResNet101, uzyskując AUC równy 0,938 oraz dokładność na poziomie 0,768 w zbiorze testowym.
Model ten został wykorzystany w implementacji aplikacji mobilnej.
Skuteczność aplikacji oceniono w ramach wstępnych testów empirycznych przeprowadzonych na niewielkiej próbie łagodnych zmian skórnych, których charakter został potwierdzony przez dermatologa.
Wyniki te potwierdziły poprawność działania systemu w warunkach zbliżonych do rzeczywistego użytkowania.

Kolejnym etapem badań powinno być rozszerzenie testów empirycznych o rzadsze i złośliwe zmiany skórne, przeprowadzane pod nadzorem dermatologa.
W przyszłości aplikację można również rozszerzyć o obsługę innych systemów operacyjnych, ponieważ obecnie funkcjonuje wyłącznie na urządzeniach z systemem Android.
Istotnym kierunkiem rozwoju jest także implementacja funkcjonalności umożliwiającej monitorowanie zmian w czasie, co stanowi ważny element diagnostyki dermatologicznej, a jednocześnie nie zawsze jest możliwe w warunkach klinicznych.

Ograniczeniem obecnej wersji systemu jest wykorzystywanie wyłącznie danych obrazowych.
Uwzględnienie dodatkowych informacji, takich jak wiek, płeć, lokalizacja zmiany czy obecność bólu lub świądu, mogłoby znacząco zwiększyć skuteczność diagnostyczną, szczególnie w przypadku raka kolczystokomórkowego.
Kolejną wadą jest brak możliwości porównania danej zmiany z pozostałymi znamionami pacjenta, co uniemożliwia zastosowanie zasady "`brzydkiego kaczątka'', powszechnie wykorzystywanej przez dermatologów.

Mimo tych ograniczeń, zaprezentowany system wykazuje znaczący potencjał w obszarze diagnostyki cyfrowej.
Dostępność narzędzia na urządzeniach mobilnych może realnie zwiększać świadomość pacjentów dotyczącą wyglądu potencjalnie złośliwych zmian skórnych.
Tworzenie tego typu prototypów jest szczególnie istotne w kontekście dynamicznego rozwoju technologii sztucznej inteligencji oraz kształtujących się ram regulacyjnych.
Systemy takie mogą w przyszłości przyczynić się do poprawy wczesnej diagnostyki nowotworów skóry na całym świecie, zwłaszcza w regionach o ograniczonym dostępie do specjalistycznej opieki dermatologicznej.